\documentclass[12pt,a4paper]{article}
\usepackage[T1]{fontenc}
\usepackage[utf8]{inputenc}
\usepackage{lmodern}
\usepackage[a4paper,margin=1in]{geometry}
\usepackage{amsmath,amssymb,amsthm}
\usepackage{microtype}
\usepackage{enumitem}
\usepackage{setspace}
\usepackage{mathtools}
\allowdisplaybreaks
\setlength{\parindent}{1.5em}
\setlength{\parskip}{0pt}
\raggedbottom

\newtheorem{theorem}{Theorem}
\newtheorem{lemma}{Lemma}

\begin{document}

\begin{titlepage}
\thispagestyle{empty}
\noindent\makebox[\textwidth][r]{\large SERIE A --- MATHEMATIK}

\vspace*{0.20\textheight}

\begin{center}
{\LARGE Expanding graphs and Ramsey numbers\par}
\vspace{1.2em}
{\large Stephan Brandt$^{\dagger *}$\par}
\vspace{2em}
{\Large Preprint No.\ A 96--24\par}
\vspace{0.3em}
{\Large December 1996\par}
\vfill
{\large\bfseries Abstract\par}
\vspace{0.8em}
\begin{minipage}{0.80\textwidth}
The generalized Ramsey number $r(G;H)$ is investigated for $H$ being a large order
graph of bounded maximum degree. Giving a negative answer to a number of con-
jectures in generalized Ramsey theory it will be shown that for every nonbipartite
graph $G$ there is a function $h(G,d)$ tending to infinity with $d$, such that the Ramsey
number $r(G;H)>h(G,d)|H|$ for almost every $d$-regular graph $H$.
\end{minipage}
\end{center}

\vfill

\noindent$^{\dagger}$ FB Mathematik und Informatik, WE 2,\\
Freie Universit\"at Berlin,\\
Arnimallee 3, 14195 Berlin, Germany\\
e-mail: brandt@math.fu-berlin.de\\
$^{*}$Partly supported by Deutsche Forschungsgemeinschaft, Grant We 1265
\end{titlepage}

\setcounter{page}{2}

\begin{center}
{\Large\bfseries Dedication}
\end{center}

This paper is dedicated to the memory of Paul Erd\H{o}s, who died in Warsaw on
September 20, 1996 while attending a conference. I communicated with him about
these results on the day before his death, where he suggested to estimate the function
$h$, which is performed in Theorem 2 below.

\section{Introduction}

For two graphs $G$ and $H$ the Ramsey number $r(G;H)$ is the smallest integer $r$ such
that every graph $F$ of order $p\ge r$ contains $G$ as a subgraph or its complement $\overline{F}$
contains $H$ as a subgraph. Chv\'atal, R\"odl, Szemer\'edi and Trotter [18] proved that
the Ramsey number $r(H;H)$ is linear in the order of $H$ for graphs of bounded degree.
More precisely, for every $\Delta$ there is a constant $c(\Delta)$, such that $r(H;H)\le c(\Delta)|H|$ for
every graph $H$ with maximum degree at most $\Delta$. Certainly, $c(\Delta)$ is an exponentially
growing function of $\Delta$ (though the growth might be slower if $|H|$ is required to be
much larger than $\Delta$).

Here we investigate the growth of Ramsey numbers $r(G;H)$, where $G$ is a fixed
graph and $H$ is a sequence of graphs of bounded maximum degree. It is easy
to derive from the result of Chv\'atal et. al., that $r(G;H)\le c(\Delta(G)+\Delta)|H|$ for
every graph $H$ with maximum degree at most $\Delta$. The objective of this paper is
to determine lower bounds for $r(G;H)$. Using the expansion properties of random
regular graphs, we will show that for every nonbipartite graph $G$ the Ramsey number
$r(G;H)\ge h(G,d)|H|$ for almost every (with probability tending to 1) $d$-regular
graph $H$, where $h(G,d)$ tends to infinity with $d$. This contradicts a conjecture of
Burr [11, Conj. 4.1], saying that $r(G;H)<\chi(G)|H|$ for almost every (all but finitely
many) connected graph $H$ of maximum degree at most $\Delta$ ($\chi$ being the chromatic
number), and several other conjectures as well.

This research was motivated by the wealth of experimental data on Ramsey
numbers $r(K_3;H)$ [6, 7, 8] for (relatively) small order graphs $H$, where the author
first saw well-expanding graphs $H$ having largest Ramsey numbers among undense
graphs.

\section{History and main results}

Let $H$ be a connected graph of order $n$. Generalizing an observation of Chv\'atal and
Harary [17], Burr [10] established the lower bound of
\begin{equation}
  r(G;H)\ge (\chi(G)-1)(n-1)+s(G)
\end{equation}
for any graph $G$, where $\chi(G)$ denotes the chromatic number and $s(G)$ the chromatic
surplus of $G$, i.e. the minimum cardinality of a color class taken over all proper $\chi(G)$
colorings of $G$. Burr [10] defined $H$ to be $G$-good, if (1) holds with equality.

Which graphs are $G$-good? Obviously, every graph is $K_2$-good. Furthermore,
Burr [10] proved that if $H$ is $K_m$-good then it is also $K_{m'}$-good whenever $2\le m'\le m$.

Two results, dating back to the times before the goodness concept was defined,
are saying that the most undense connected graphs are $K_m$-good, namely trees
(Chv\'atal [16]) and cycles of length at least $m^2-2$ (Bondy and Erd\H{o}s [5]).

Three conjectures have been posed describing the expected limits for $G$-goodness
(see [11, Section 4]).

\medskip
\noindent\textbf{Conjecture 1 (Burr [10])} If $G$ is any fixed graph and $c$ is any fixed number, then
any connected graph of sufficiently large order with maximum degree $\Delta\le c$ is $G$-
good.

\medskip
\noindent\textbf{Conjecture 2 (Burr and Erd\H{o}s [12])} If $G=K_m$ and $c$ is any fixed number, then
any connected graph $H$ of sufficiently large order where every subgraph $H'$ satisfies
$|E(H')|/|H'|<c$ is $G$-good.

\medskip
\noindent\textbf{Conjecture 3 (Burr [11])} If $G=K_3$ and $c$ is any fixed number, then any con-
nected graph of sufficiently large order $n$ of size at most $cn$ is $G$-good.

Two further conjectures by Burr and Faudree [15, Conj. 1.1 and 2.1] generalize
Conjecture 2 (for which a reward of 25\$ was offered by the authors) to larger classes
of graphs than just $K_m$.

Conjecture 1 was shown to hold if $G$ is bipartite by Burr, Erd\H{o}s, Faudree,
Rousseau and Schelp [14]. Here we will show that it is false in every other case.
Even more, if $G$ is not bipartite then Conjecture 1 is false for almost every $d$-regular
graph $H$, if $d$ is sufficiently large with respect to $G$. In fact, the Ramsey number
$r(G;H)$ exceeds $\gamma n$ for any constant $\gamma$, if $d$ is sufficiently large with respect to $G$
and $\gamma$ for almost every $d$-regular graph. So, in particular, all the other conjectures
are false as well. The term \emph{almost every} means ``with probability $1-o(1)$'' unless
otherwise stated.

\begin{theorem}
For every nonbipartite graph $G$ there is a function $h(G,d)$ tending to
infinity as $d\to\infty$, such that $r(G;H)>h(G,d)n$ for almost every $d$-regular graph
$H$ of order $n$.
\end{theorem}

In order to prove Theorem 1 it suffices to prove the statement for odd cy-
cles $C_{2k+1}$ ($k\ge 1$), since every nonbipartite graph contains an odd cycle, and
$r(G;H)\ge r(G';H)$ for every subgraph $G'$ of $G$. We can give explicit lower bounds
for $h(C_{2k+1},d)$. In fact, it turns out that we can choose $h$ to be a reasonably fast
growing function of $d$.

\begin{theorem}
For odd cycles, Theorem 1 holds with
\[
  h(C_3,d)=\Omega(\sqrt{d}/\log d)
\]
and
\[
  h(C_{2k+1},d)=\Omega((d/\log d)^{1/(4k+2)})
\]
for every $k\ge 2$.
\end{theorem}

Moreover, Burr, Erd\H{o}s, Faudree, Rousseau and Schelp [13] asked for the functions
$f(n,m)$ and $g(n,m)$, where $f(n,m)$ is the largest integer such that every connected
graph of order $n$ and size at most $f(n,m)$ is $K_m$-good, and $g(n,m)$ is the largest
integer such that some connected graph of order $n$ and size $f(n,m)$ is $K_m$-good.

While $g(n,m)$ is for every fixed $m$ a superlinear function in $n$ [13], the situation is
different for $f(n,m)$. For $m\ge 4$, Burr et.al. [13] proved that $f(n,m)=n+o(n)$. For
$m=3$ they proved that $f(n,3)/n>17/15$. We will show that $f(n,3)/n<84$, and,
in fact, with a different approach and a more refined analysis which is not presented
here we can show that $f(n,3)/n<11.75$. However, this is likely to be still far from
the optimum. The author would expect $2<f(n,3)/n<6$, say, if $n$ is sufficiently
large. Computer experiments [7] seem to indicate that at least $f(n,3)/n>3/2$
holds for larger $n$.

The method that we use is as follows: We employ the expansion properties
of random regular graphs to justify that these graphs cannot be contained in the
complement of lexicographic products $F[K_s]$, where $F$ is a certain graph with larger
odd girth than $G$ and small independence number.

\section{Concepts and notation}

The graph theoretical notation is standard as far as possible; for notation and con-
cepts not defined here the reader is referred to introductory literature, e.g. [24].

The \emph{lexicographic product} (composition) $G[H]$ of two graphs $G$ and $H$ has vertex
set $V(G)\times V(H)$ and two vertices $(v_i,w_i)$ and $(v_j,w_j)$ are joined by an edge if and
only if (1) $v_iv_j\in E(G)$ or (2) $v_i=v_j$ and $w_iw_j\in E(H)$. In our case we will consider
only the case where $H$ is a complete graph or its complement.

The \emph{girth} of a graph is the length of a shortest cycle and the \emph{odd girth} is the
length of a shortest odd cycle. Note that the odd girth of a graph $G$ is inherited by
the lexicographic product $G[\bar{K}_r]$.

The \emph{clique number} $\omega(G)$ of a graph $G$ is the maximum number of vertices which
are pairwise joined by an edge and the \emph{independence number} $\alpha(G)$ is the maximum
number of vertices which are pairwise not joined by an edge. The \emph{independence ratio}
is $\alpha(G)/|G|$, where $|G|$ denotes the order of $G$.

A \emph{subgraph} of a graph $G$ is a graph $G'$ with $V(G')\subseteq V(G)$ and $E(G')\subseteq E(G)\cap
\binom{V(G')}{2}$, where the edges of the undirected, simple graph are considered as two-
element subsets of the vertex set, and $\binom{V}{2}$ denotes the set of all two-element subsets
of $V$. An \emph{induced subgraph} is a subgraph $G'$ satisfying $E(G')=E(G)\cap\binom{V(G')}{2}$.

\section{Expanding graphs}

Typical graphs tend to have the property that the edges are homogeneously dis-
tributed over the graph (even though most infinite sequences of graphs that we know
do not have this property). This means, in particular, that for any two reasonably
large disjoint subsets $U,W$ of vertices the number of edges joining them does not
deviate a lot from the expected number of edges (which is $|U||W||E(H)|/\binom{|V(H)|}{2}$).
A necessary condition for $H$ being a subgraph of a graph $F$ is the following:

\begin{quote}
$F$ has an induced subgraph $F'$ of order $|V(H)|$, such that for every pair
$u,w$ the minimum number of edges joining two sets of cardinalities $u$
and $w$ is at least as large in $F'$ as in $H$.
\end{quote}

This is the key obstruction that will be used to prove that the Ramsey number
$r(C_{2k+1};H)$ is large for well-expanding graphs $H$.

Which families of graphs with linearly many edges will be suitable for having
a large Ramsey number (vs.\ $C_{2k+1}$), due to this reason? A random regular graph
will almost surely do. Moreover, as shown by Robinson and Wormald [23], for
$d\ge 3$ a random $d$-regular graph is almost surely hamiltonian, so, in particular, it is
connected.

A suitable model for random regular graphs was introduced by Bollob\'as [2]
(see also [3]). To construct a $d$-regular graph with $n$ vertices ($d$ times $n$ even)
take a set $\{1,2,\ldots,n\}$ of vertices with $d$ half-edges sticking out of each vertex.
Choose a random pairing of the half-edges to edges. Such a pairing will be called a
\emph{configuration}. The result will be a $d$-regular multigraph, possibly containing loops or
multiple edges. However, for fixed $d$, a positive fraction of the graphs will be simple
graphs [2]. So any statement, which holds for almost every $d$-regular multigraph,
will hold for almost every $d$-regular simple graph as well. Note that every labelled
$d$-regular simple graph appears with the same probability in this model. Moreover,
the number of configurations with $2x$ half edges is
\[
  N(2x)=\frac{(2x)!}{x!2^x}.
\]
From Stirling's formula, we get $(1-o(1))\sqrt{2}\left(\frac{2x}{e}\right)^x<N(2x)<\sqrt{2}\left(\frac{2x}{e}\right)^x$.

Bollob\'as [4] proved that almost every $d$-regular graph for $d\ge 3$ has large isoperi-
metric number. The next result shows that such graphs are good expanders in a
stronger sense. To avoid tedious case analysis we concentrate on the range of $c$ that
we will need in our proofs and we make no attempt to optimize the bound for $d$.

\begin{theorem}
Suppose $0<c\le 1/15$ and $d$ is an integer satisfying $d>2(1-\ln c)/c(1-5c)$. Let $H$ be a random $d$-regular multigraph of order $n$ ($dn$ even).
Then almost surely (as $n\to\infty$) there is an edge between any pair of disjoint vertex
sets $U,W$ satisfying $|U|=|W|\ge cn$.
\end{theorem}

\begin{proof}
Suppose $|U|=|W|=m=\lceil cn\rceil=\eta n$, so $\eta\to c$ as $n\to\infty$. Consider the
configurations where exactly $k$ edges join vertices of $U$. Suppose that no edge joins
$U$ to $W$. Then all ends of the remaining $dm-2k$ half-edges emanating from $U$ have
the other end in $V(H)\setminus(U\cup W)$, and for a chosen subset of $dm-2k$ half-edges
emanating from $V(H)\setminus(U\cup W)$ we have $(dm-2k)!$ possibilities to pair them with
the half-edges of $U$. Finally, we can pair the remaining half-edges in any fashion.
Set $2k/d=\gamma n$ and let $X_{U,W}$ denote the number of edges joining $U$ to $W$. Then the
probability
\[
  P[X_{U,W}=0]=\sum_{0\le k\le dm/2}
  \frac{\binom{dm}{2k}N(2k)\binom{d(n-2m)}{dm-2k}(dm-2k)!N(d(n-2m)+2k)}{N(dn)}.
\]
We can estimate $\binom{dm}{2k}\le \left(\frac{e\eta}{\gamma}\right)^{\gamma dn}$,
$N(2k)\le \sqrt{2}\left(\frac{\gamma dn}{e}\right)^{\gamma dn/2}$,
$\binom{d(n-2m)}{dm-2k}(dm-2k)!\le n\bigl(dn(1-5\eta/2+\gamma/2)\bigr)^{(\eta-\gamma)dn}$
by pairing the terms around the middle, and
\begin{align*}
  \frac{N(d(n-2m)+2k)}{N(dn)}
  &<(1+o(1))\left(\frac{dn(1-2\eta+\gamma)}{e}\right)^{dn(1-2\eta+\gamma)/2}
    \left(\frac{e}{dn}\right)^{dn/2} \\
  &= (1+o(1))(1-2\eta+\gamma)^{dn/2}\left(\frac{e}{dn(1-2\eta+\gamma)}\right)^{(\eta-\gamma/2)dn} \\
  &\le (1+o(1))\left(e^{-(2\eta-\gamma)^2/4}\bigl(dn(1-2\eta+\gamma)\bigr)^{-(\eta-\gamma/2)}\right)^{dn},
\end{align*}
where the last term follows from the inequality $1-x\le e^{-x-x^2/2}$ for $0\le x<1$.
Writing
\[
  \Psi(\eta,\gamma)=\left(\frac{e\eta^2}{\gamma(1-2\eta)}\right)^\gamma
\]
we get that $\sqrt{2}(\Psi(\eta,\gamma))^{dn/2}>\binom{dm}{2k}N(2k)/((1-2\eta+\gamma)dn)^{\gamma dn/2}$, so
\begin{align*}
  P[X_{U,W}=0]
  &< O(n^2)\max_{0\le \gamma\le \eta}
     \left[(\Psi(\eta,\gamma))^{dn/2}
     \left(\frac{1-5\eta/2+\gamma/2}{1-2\eta+\gamma}\right)^{(\eta-\gamma)dn}
     e^{-dn(2\eta-\gamma)^2/4}\right] \\
  &< O(n^2)\max_{0\le \gamma\le \eta}
     \left(\Psi(\eta,\gamma)(1-\eta/2-\gamma/2)^{2(\eta-\gamma)}e^{-(2\eta-\gamma)^2/2}\right)^{dn/2} \\
  &< O(n^2)\max_{0\le \gamma\le \eta}
     \left(\Psi(\eta,\gamma)e^{-(\eta^2-\gamma^2)-(2\eta-\gamma)^2/2}\right)^{dn/2}.
\end{align*}
The function $\Psi(\eta,\gamma)$ is maximal if $\gamma=\eta^2/(1-2\eta)$, whence $\Psi(\eta,\gamma)\le e^{\eta^2/(1-2\eta)}$.
Now suppose that $n$ is sufficiently large. If $\gamma\le 7\eta^2/2$ then using $\eta<(1+o(1))/15$
we get
\[
  P[X_{U,W}=0]<e^{((1-2\eta)^{-1}-(3-7\eta-49\eta/120))\eta^2dn/2}<e^{-(1-5\eta)\eta^2dn}.
\]
If $\gamma>7\eta^2/2$ then $\Psi(\eta,\gamma)$ has a negative contribution in the exponent and $P[X_{U,W}=
0]<e^{-(1-5\eta)\eta^2dn}$ holds as well. On the other hand, there are
\[
  \frac{\binom{n}{m}\binom{n-m}{m}}{2}=\frac{\binom{n}{2m}\binom{2m}{m}}{2}<\left(\frac{e}{2\eta}\right)^{2\eta n}2^{2\eta n}=e^{(1-\ln\eta)2\eta n}
\]
pairs of disjoint sets of cardinalities $m$. So the probability that there is a pair of
disjoint sets of cardinalities $m$, which are not joined by an edge is at most
\[
  \left(e^{2(1-\ln\eta)-d(1-5\eta)\eta}\right)^{\eta n}.
\]
Hence, if $d>2(1-\ln c)/c(1-5c)$ then in a random $d$-regular multigraph almost
surely every pair of disjoint subsets of cardinalities at least $cn$ is joined by an edge.
\end{proof}

This is only an existence proof, and given a random regular graph it may be diffi-
cult to decide, whether the graph has this property. Anyway, we can at least justify
that a regular graph has this property based on the eigenvalues of the adjacency ma-
trix (see [1, p.~122]). Broder and Shamir [9] (cf. the improvement of Friedman [20])
proved in a slightly different and more restricted model of random regular graphs,
that the justification via eigenvalues is possible for almost all $d$-regular graphs if
$d$ is sufficiently large and even. Moreover, explicit constructions of infinite fam-
ilies of regular graphs with constant degree are known which are well-expanding
in the sense of Theorem 3, e.g. the non-bipartite Ramanujan graphs constructed
by Lubotzky, Phillips and Sarnak [22]. Note that the bounds for $d$ obtained via
eigenvalue estimates are weaker than those computed directly in Theorem 3.

\section{Lexicographic products as Ramsey graphs}

In the previous section we analyzed properties of well-expanding graphs with re-
gard to our needs. In order to show that for such a graph $H$ the Ramsey number
$r(C_{2k+1};H)$ is large, one has to find a graph of large order compared to $H$, which
does not contain $C_{2k+1}$ and whose complement does not contain $H$. This graph
will be the lexicographic product of a graph with small independence number and
odd girth exceeding $2k+1$ with an edgeless graph. Note that $F[\bar{K}_r]$ has the same
odd girth as $F$, while for $r\ge 2$ it has many short even cycles. Moreover, the
independence ratio $\alpha(F[\bar{K}_r])/|F[\bar{K}_r]|=\alpha(F)/|F|$.

Certainly, $C_{2k+1}$ is not a subgraph of $F[\bar{K}_r]$. Turning to the complement, we
will see that $\overline{F[\bar{K}_r]}=\overline{F}[K_r]$ is not a well expanding graph in the required sense,
since every subgraph of order $n=|F[\bar{K}_r]|/h$ contains two disjoint vertex sets of
cardinalities $\lceil cn\rceil$ which are not joined by an edge, $c$ being a constant only depending
on $F$. So any sufficiently well-expanding graph $H$ of order $n$ cannot be a subgraph of
$\overline{F}[K_r]$. We simulate subgraphs of lexicographic products $\overline{F}[K_r]$ by assigning weights
to the vertices $v_i$ of $F$, representing the ratio of the vertices in $V_i$ belonging to the
subgraph.

\begin{lemma}
Let $F$ be a graph of order $p$ with clique number $\omega$. Assign to every vertex
$v$ of $F$ a weight $w_v$ with $0\le w_v\le 1$. If $t=\sum_{v\in V}w_v>\omega$, then $F$ contains two
non-adjacent vertices $x_1,x_2$ with $w_{x_i}\ge (t-\omega)/(p-\omega)$.
\end{lemma}

\begin{proof}
Let $V'$ be the set of vertices $v$ with $w_v\ge (t-\omega)/(p-\omega)$. Since $\omega+(p-
\omega)(t-\omega)/(p-\omega)=t$ we obtain $|V'|>\omega$. So $V'$ contains two non-adjacent vertices.
\end{proof}

Before treating the general case, we first prove that $f(n,3)/n<84$ by showing
that $r(K_3;H)>2|H|$ for almost every $d$-regular graph with $d\ge 168$. Let $r=\lceil 2n/5\rceil$
and take $F=C_5$. Setting $t=|H|/r=5(1-o(1))/2$ and assigning weights of total
weight $t$ to the vertices of $C_5$, we get from Lemma 1 that $C_5$ contains two non-
adjacent vertices of weight $w_{x_i}\ge (1-o(1))/6$. So every subgraph $F'$ of $\overline{C_5[\bar{K}_r]}=$
$C_5[K_r]$ of order $|F'|=tr=n$ contains two disjoint subsets, each of cardinality
$|F'|w_{x_i}/t=(1-o(1))n/15$, which are not joined by an edge. Employing Theorem 3
with $c=(1-o(1))/15$ we get that, as $n\to\infty$, for almost every $d$-regular graph $H$
with $d\ge 168$ and order $n$ every pair of disjoint subsets of cardinalities $(1-o(1))n/15$
is joined by an edge. So $H$ cannot be a subgraph of any subgraph of $C_5[K_r]$ with the
same order as $H$, and hence not be a subgraph of $C_5[K_r]$. So the Ramsey number
$r(K_3;H)>2n$ and hence $f(n,3)/n<84$.

It should be mentioned that $C_5[K_r]$ is, in fact, Ramsey graph for the Ramsey
number $r(K_3;H)$ of certain undense graphs $H$ which are well-expanding. E.g. let
$H$ be the (unique) $r(K_3,K_5)$-Ramsey graph itself which is 4-regular of order 13 and
well expanding. Using the program \texttt{mtf} due to Brandt, Brinkmann and Harmuth [6]
we can compute that $r(K_3;H)=26$, and $C_5[K_5]$ is a Ramsey graph. This shows,
by the way, that $f(13,3)/13<2$.

The same approach can be used to show that $r(C_{2k+1};H)>h(C_{2k+1},d)|H|$,
where $h(C_{2k+1},d)\to\infty$ as $d\to\infty$. We will make use of the existence of families
of graphs $F_i$ with the following properties: $\alpha(F_i)/|F_i|$ tends to 0 and $F_i[\bar{K}_r]$ does
not contain $C_{2k+1}$ for any $r\ge 1$. We can simply choose $F_i$ to be a family of graphs
with girth exceeding $2k+1$, since the odd girth of $F_i[\bar{K}_r]$ remains unchanged. The
existence of families with arbitrarily large girth and arbitrarily small independence
ratio was proven by Erd\H{o}s [19] by probabilistic means, in order to show that there
are graphs with arbitrarily large girth and chromatic number.

\begin{proof}[Proof of Theorem 2]
Let $\beta(k,p)$ be the minimum independence number taken over all graphs of order
$p$ with odd girth more than $2k+1$. By Kim's result [21, Theorem 1.1] we have
$\beta(3,p)=O(\sqrt{p\log p})$ (note that the inequality in [21, Theorem 1.1] is misprinted)
and by Erd\H{o}s' result [19] we have $\beta(k,p)=O(p^{1-1/(4k+2)})$.

Let $d$ be an integer which is sufficiently large with respect to $k$ and set $c=
3\ln d/d$. By Theorem 3 in a $d'$-regular graph $H$ of order $n$ with $d'\ge d$ almost surely
every pair of disjoint vertex sets $U,W$ of cardinality $|U|,|W|\ge cn$ are joined by an
edge since $d'\ge d>2(1-\ln c)/(1-5c)c$.

For $p=\lceil 1/2c\rceil$ there exists a graph $F$ of odd girth larger than $2k+1$ and order
$p$ which has independence number $\beta(k,p)$. Now set $r=\lceil n/2\beta(k,p)\rceil$ and choose
$0<\varepsilon<\beta(k,p)/2p$. If $n$ is sufficiently large then $n/r\ge 2(1-\varepsilon)\beta(k,p)$. Assume
that $H$ is a subgraph of $\overline{F[\bar{K}_r]}$, i.e. there is an injection $\sigma:V(H)\to V(\overline{F[\bar{K}_r]})$ such
that $xy\in E(H)$ implies $\sigma(x)\sigma(y)\in E(\overline{F[\bar{K}_r]})$. Let $k_i$ be the number of vertices
from $H$ with $\sigma(x)=(v_i,\cdot)$. Assign to every vertex $v_i$ of $\overline{F}$ weight $w_{v_i}=k_i/r$. So
the total weight is $n/r\ge 2(1-\varepsilon)\beta(k,p)$. By Lemma 1 the graph $\overline{F}$ contains two
non-adjacent vertices $x_1,x_2$ of weight
\[
  w_{x_i}\ge \frac{2(1-\varepsilon)\beta(k,p)-\omega}{p-\omega}\ge \frac{(1-2\varepsilon)\beta(k,p)}{p-\beta(k,p)}>
  \frac{\beta(k,p)}{p}.
\]
So $H$ contains two vertex sets $U,W$ with $|U|,|W|\ge \beta(k,p)r/p\ge n/2p\ge cn$ which
are not joined by an edge.

Hence a random $d'$-regular graph $H$ is almost surely not a subgraph of $\overline{F[\bar{K}_r]}$,
so, in particular $r(C_{2k+1};H)>h(C_{2k+1},d)|H|$, where $h(C_{2k+1},d)=pr/|H|\ge
p/2\beta(k,p)$. Since $p=\Theta(d/\log d)$ we get $h(C_3,d)=\Omega(\sqrt{d}/\log d)$ and $h(C_{2k+1},d)=
\Omega((d/\log d)^{1/(4k+2)})$ for $k\ge 2$.
\end{proof}

The property that every pair of small disjoint subsets is joined by an edge is
a strong expansion property, which is not satisfied e.g. by well-expanding bipartite
graphs. It would be interesting to know whether weaker expansion properties can
already force a sequence of undense graphs $H$ to have large Ramsey number $r(G;H)$,
perhaps restricted to graphs $G$ in a subclass of nonbipartite graphs.

\medskip
\noindent\textbf{Acknowledgement}

I thank Alan Frieze and Jean-Pierre Tillich for their helpful remarks, and Hans
Mielke for pointing out a mistake in an earlier version.

\begin{thebibliography}{99}

\bibitem{1}
N. Alon and J. H. Spencer, \emph{The probabilistic method}, Wiley, New York,
1992.

\bibitem{2}
B. Bollob\'as, A probabilistic proof of an asymptotic formula for the number
of labelled regular graphs, \emph{Europ. J. Comb.} \textbf{1} (1980), 34--38.

\bibitem{3}
B. Bollob\'as, \emph{Random Graphs}, Academic Press, London, 1985.

\bibitem{4}
B. Bollob\'as, The isoperimetric number of random regular graphs, \emph{Europ. J.}
\emph{Comb.} \textbf{9} (1988), 241--244.

\bibitem{5}
J. A. Bondy and P. Erd\H{o}s, Ramsey numbers for cycles in graphs, \emph{J. Comb.}
\emph{Theory Ser. B} \textbf{14} (1974), 46--54.

\bibitem{6}
S. Brandt, G. Brinkmann, T. Harmuth, The generation of maximal
triangle-free graphs, in preparation.

\bibitem{7}
S. Brandt, G. Brinkmann, T. Harmuth, All Ramsey numbers $r(K_3,G)$
for connected graphs of order 9, in preparation.

\bibitem{8}
G. Brinkmann, All Ramsey numbers $r(K_3,G)$ for connected graphs of order
7 and 8, submitted.

\bibitem{9}
A. Broder and E. Shamir, On the second eigenvalue of random regular
graphs, In: 28th Symp. on Foundations of Comp. Sci., IEEE (1987), pp.~286--
294.

\bibitem{10}
S. A. Burr, Ramsey numbers involving graphs with long suspended paths, \emph{J.}
\emph{London Math. Soc.} (2) \textbf{24} (1981), 11--20.

\bibitem{11}
S. A. Burr, What can we hope to accomplish in generalized Ramsey theory?,
\emph{Discrete Math.} \textbf{67} (1987), 215--225.

\bibitem{12}
S. A. Burr and P. Erd\H{o}s, Generalizations of a Ramsey-theoretic result of
Chv\'atal, \emph{J. Graph Theory} \textbf{7} (1983), 39--51.

\bibitem{13}
S. A. Burr, P. Erd\H{o}s, R. J. Faudree, C. C. Rousseau and R. H.
Schelp, An extremal problem in generalized Ramsey theory, \emph{Ars Comb.} \textbf{10}
(1980), 193--203.

\bibitem{14}
S. A. Burr, P. Erd\H{o}s, R. J. Faudree, C. C. Rousseau and R. H.
Schelp, The Ramsey number for the pair complete bipartite graph---graph of
limited degree, in: \emph{Graph Theory with Applications to Algorithms and Com-}
\emph{puter Science} (Alavi et. al., eds.) Wiley, 1985, 163--174.

\bibitem{15}
S. A. Burr and R. J. Faudree, On graphs $G$ for which all large trees are
$G$-good, \emph{Graphs \& Comb.} \textbf{9} (1993), 305--313.

\bibitem{16}
V. Chv\'atal, Tree---complete graph Ramsey numbers, \emph{J. Graph Theory} \textbf{1}
(1977), 93.

\bibitem{17}
V. Chv\'atal and F. Harary, Generalized Ramsey theory for graphs III,
Small off-diagonal numbers, \emph{Pacific J. Math.} \textbf{41} (1972), 335--345.

\bibitem{18}
V. Chv\'atal, V. R\"odl, E. Szemer\'edi and W. T. Trotter, The Ramsey
number of graphs with bounded degree, \emph{J. Comb. Theory Ser. B} \textbf{34} (1983),
239--243.

\bibitem{19}
P. Erd\H{o}s, Graph theory and probability, \emph{Canad. J. Math.} \textbf{11} (1959), 34--38.

\bibitem{20}
J. Friedman, On the second eigenvalue and random walks in random $d$-regular
graphs, \emph{Combinatorica} \textbf{11} (1991), 331--362.

\bibitem{21}
J. H. Kim, The Ramsey number $R(3,t)$ has order of magnitude $t^2/\log t$, \emph{Ran-}
\emph{dom Structures \& Algorithms} \textbf{7} (1995), 173--207.

\bibitem{22}
A. Lubotzky, R. Phillips and P. Sarnak, Ramanujan graphs, \emph{Combina-}
\emph{torica} \textbf{8} (1988), 261--277.

\bibitem{23}
R. W. Robinson and N. C. Wormald, Almost all regular graphs are Hamil-
tonian, \emph{Random Structures and Algorithms} \textbf{5} (1994), 363--374.

\bibitem{24}
D. B. West, \emph{Introduction to graph theory}, Prentice Hall, 1996.

\end{thebibliography}

\end{document}
